\subsection{Knowledge representation}\label{ontology}

We represent the resource as a modular RDF knowledge base. The
terminology module (TBox) defines classes and relations in OWL
\cite{owl2}; the domain modules (ABoxes) describe sites, physical
specimens, processes and data products. Shared identifiers connect
these modules (Table~\ref{tab:modules}). SIO supplies the upper-level
classes and relations \cite{sio}. The terminology declares 333
classes (297 project classes and 36 referenced classes), 20 object
properties and 35 datatype properties. These counts apply to
declarations in the terminology module.

\subsubsection{Measurement and provenance}

We separate a measured quality from the information that quantifies
it. A temperature quality belongs to a site; a measurement process
produces a value for that quality. The value contains the numerical
reading and, where documented, its unit. Counts, dimensional
measurements and measurements with undocumented units use the
appropriate datatype and unit constraints in their assigned shapes.

Listing~\ref{lst:ttl_temp} illustrates this distinction for a field
temperature. The process belongs to a site visit, the visit targets
the site, and the measured quality is an attribute of that site.
The process-to-visit path supplies the observation context. SIO's
value-to-quality relation and its inverse support traversal in both
directions.

\begin{lstlisting}[caption={Field temperature and its visit provenance.},label={lst:ttl_temp},basicstyle=\ttfamily\scriptsize,frame=single]
@prefix rak: <https://rubalkhali.science/kb/RAK_> .
@prefix sio: <http://semanticscience.org/resource/SIO_> .
@prefix pato: <http://purl.obolibrary.org/obo/PATO_> .
@prefix uo: <http://purl.obolibrary.org/obo/UO_> .
@prefix xsd: <http://www.w3.org/2001/XMLSchema#> .

rak:P000001 a rak:0000006 ;
  sio:000068 rak:3000001 ;
  sio:000229 rak:4000001 .
rak:3000001 a rak:0000003 ;
  sio:000291 rak:1000001 .
rak:4000001 a rak:0000010 ;
  sio:000221 uo:0000027 ;
  rak:2000003 "20.7"^^xsd:double ;
  sio:000215 rak:5000001 .
rak:5000001 a pato:0000146 ;
  sio:000011 rak:1000001 ;
  sio:000216 rak:4000001 .
\end{lstlisting}

Sampling processes output material-sample collections. Collection
membership links each physical replicate to the sampling event;
the specimen's derivation edge also links it directly to its site.
A protocol is a reusable information specification, and a protocol
execution is a process. Extraction consumes soil and produces DNA;
library preparation consumes DNA and produces an amplicon library;
sequencing consumes the library and produces a digital FASTQ dataset.
Quality-control processes consume FASTQ datasets and produce
direction-specific metrics. Suppl. Sections~S1.2--S1.7 give the
corresponding formal patterns and graph examples.

Site annotations use existential restrictions to relate a site to
an ENVO biome or environmental feature. Such a restriction specifies
a filler of the named class. The source annotation ledger provides
the recorded assignment and its provenance.

Laboratory controls are occurrence-specific materials with
processing roles. A control-processing execution consumes the control,
realizes its role and belongs to a documented processing group.
Trips~4--5 use trip-specific extraction-day batches. Trips~1--3 use
per-kit controls, and Trips~1--2 include shared extraction days.
Occurrence-specific identifiers and the identity ledger distinguish
reused control labels. Confirmed, provisional, unresolved and
inapplicable dispositions retain their supporting evidence. The
D6300 whole-cell and D6322 purified-DNA standards enter the process
chain at extraction and PCR, respectively.

\subsubsection{Taxonomic and chemical observations}

For each profile, the taxonomy generator aggregates ASV counts at
seven ranks from domain to species, retaining the rank-truncated
classifier lineage and its validated NCBI or project identifier.
Each workflow produces fourteen datasets: read counts and relative
read abundance for each rank. The relative value divides the aggregated
count by the profile's total reads; explicit unclassified lineages
retain their contribution to that denominator. A taxon-rank entry adds a value and a
quality linked to shared dataset, workflow, taxon and FASTQ objects.
These links preserve the profile context during retrieval.
The taxon IRI is used as an individual in the abundance graph and as
a class in the mapping module. OWL~2 punning gives these uses separate
semantics: observation links identify the measured taxon, while
class-level axioms supply the reference ancestry. The generator
performs the rank aggregation explicitly.

XRF processes have several concentration-value outputs. A laboratory
process takes an archived specimen as input; a field process targets
a site and retains the instrument test identifier. The chemical
ledger supplies supported ChEBI and PubChem cross-references for
reporting channels, including oxide-equivalent channels. Unit
documentation and acquisition scope are described in
Section~\ref{methods:xrf}.

Archived-soil pH processes consume and target the assayed specimen,
produce a pH value linked to its acidity quality and belong to a
reconstructed measurement session. Field-replicate substitutions
require separate physical identities and an explicit campaign,
site and compartment matching level (Section~\ref{methods:ph}).

\subsubsection{OWL and RDF implementation}\label{subsec:owl_rdf}

Groovy scripts using OWL API v5.1.20 \cite{owlapi} generate the modules.
The modular organization separates domain data from shared terminology
\cite{cuencagrau2008modular}. OpenLink Virtuoso \cite{virtuoso}
provides RDF storage and SPARQL query execution.

KG v3.0.0 serves the asserted union. A separate reasoning workflow
uses the bundled rule program, v1.5.0, to compute consequences of a
specified set of positive RDF rules. These rules expand equivalent
classes and properties, subclass and subproperty hierarchies,
instance types, domains, resource-valued ranges, inverse and symmetric
properties, and transitive properties. The schema audit generates
rules for each nonempty finite property-chain pattern in the loaded
schema, including inverse operands. Additional rules introduce and
classify value restrictions, classify existential restrictions from
existing edges and typed fillers, propagate universal restrictions
to existing resource fillers, and handle intersections, union
membership and resource enumerations. All conclusions use existing
RDF terms. The schema audit lists the loaded imports and the
constructs outside this rule set.

The asserted graph contains the deduplicated union of the loaded
modules. Rules read this graph together with a separate inferred
graph and insert only conclusions absent from both. Each update
selects a bounded batch of distinct conclusions; each rule repeats
until its next batch adds zero triples. The complete schedule repeats
until every rule adds zero triples in one full pass. This stopping
criterion defines the least fixed point of the emitted rule program.
The instance-subclass rule is partitioned by instantiated named
classes, with a separate partition for anonymous classes. The class
inventory is refreshed each pass, so types introduced later in the
schedule enter subsequent passes.

The controller preserves the rule text, partition inventories,
per-batch additions, runtime versions and input identifiers.
Query truncation, SQL errors and exhaustion of the pass limit
terminate the run with a failure status. After convergence, separate
diagnostics test explicit disjoint-class and complement clashes,
membership in \texttt{owl:Nothing}, and incompatible laboratory-control
roles. These diagnostics and the finite rule set define the scope
of operational entailment validation.

\subsubsection{Structural constraints}\label{subsec:schema_def}

ShEx shapes \cite{shex} check selected nodes for required types,
properties, cardinalities and datatypes. The sampling-site shape
requires a site type, label and geometry predicate. Its unconstrained
geometry object checks presence; validation of the WKT datatype and
coordinate syntax is a separate requirement. Additional type
assertions are accepted independently of how they were obtained.

The sample shape requires a replicate type, a label, a source
identifier and an IRI-valued derivation edge. The DNA concentration
shape checks its numeric value, nanograms-per-microlitre unit and
quality/process links. The XRF-process shape requires a protocol,
one or more values and an input or target IRI. Checking the class of
that input or target requires a corresponding target shape or
scientific-invariant check. These open shapes admit predicates
outside their declared constraints; the negative fixtures test the
specific failures reported in Technical Validation.
